\documentclass[conference]{IEEEtran}
\usepackage{cite}
\usepackage{amsmath,amssymb,amsfonts}
\usepackage{algorithmic}
\usepackage{graphicx}
\usepackage{textcomp}
\usepackage{xcolor}
\usepackage{booktabs}
\usepackage{tabularx}
\usepackage{url}

\def\BibTeX{{\rm B\kern-.05em{\sc i\kern-.025em b}\kern-.08em
    T\kern-.1667em\lower.7ex\hbox{E}\kern-.125emX}}

\begin{document}

\title{Design and Implementation of a Synthesizable 16-Bit Direct-Mapped Cache Controller Using Verilog HDL}

\author{
    \IEEEauthorblockN{Palaparthi Harsha Karthikeya}
    \IEEEauthorblockA{\textit{Roll No: 2401010313} \\
    \textit{Newton School of Technology} \\
    Rishihood University \\
    palaparthi.harshakarthikeya2024@nst.rishihood.edu.in}
    \and
    \IEEEauthorblockN{Dakarapu Hemamdhar nath}
    \IEEEauthorblockA{\textit{Roll No: 2401010137} \\
    \textit{Newton School of Technology} \\
    Rishihood University \\
    dakarapu.nath2024@nst.rishihood.edu.in}
    \and
    \IEEEauthorblockN{Akshay Y}
    \IEEEauthorblockA{\textit{Roll No: 2401020008} \\
    \textit{Newton School of Technology} \\
    Rishihood University \\
    akshay.y2024@nst.rishihood.edu.in}
}

\maketitle

% Faculty Clarification Footnote / Note
\begingroup\renewcommand\thefootnote{\textdagger}
\footnotetext{Academic Notice: Group Alpha was registered on 10/16/2026. While standard departmental templates suggest a five-member group size, our registered group strictly comprises three members: Palaparthi Harsha Karthikeya (2401010313), Dakarapu Hemamdhar nath (2401010137), and Akshay Y (2401020008). Individual contributions are explicitly detailed in Section VI.}
\endgroup

\begin{abstract}
The escalating disparity between processor execution speeds and dynamic random-access memory (DRAM) access latencies---widely recognized as the ``memory wall''---represents a fundamental bottleneck in contemporary digital computer systems. Cache memory serves as the principal hardware architectural mechanism to mitigate this latency disparity by exploiting the principles of spatial and temporal locality. This paper presents the design, register-transfer level (RTL) implementation, and cycle-accurate functional verification of a synthesizable 16-bit direct-mapped cache controller specified in Verilog-2001 (IEEE Std 1364-2001). The controller interfaces an 8-bit processing unit to a 64~KB secondary memory space through a 4-line, 4-byte/line (16-byte total data capacity) cache storage hierarchy. A deterministic six-state Moore/Mealy hybrid finite state machine (FSM) coordinates tag evaluation, line allocation, burst memory block fetches, and memory updates utilizing a Write-Through and Write-Allocate policy. Functional verification was executed using an automated, self-checking testbench comprising twenty individual verification assertions across twelve architectural test vectors, including cold misses, spatial locality hits, temporal hits, write-through coherence, and direct-mapped conflict replacements (demonstrated with addresses \texttt{0x0034} and \texttt{0x0044}). The design achieved a 100\% functional pass rate with zero cycle deadlocks.
\end{abstract}

\begin{IEEEkeywords}
Direct-mapped cache, Cache controller, Verilog HDL, Finite State Machine (FSM), Write-Through, Write-Allocate, Computer Architecture, Memory Wall.
\end{IEEEkeywords}

\section{Introduction}
Modern microprocessors operate at multi-gigahertz frequencies, capable of issuing multiple instructions per clock cycle. Conversely, primary off-chip dynamic RAM access times have improved at a significantly slower rate, leading to an increasing latency disparity quantified in computer architecture as the ``memory wall'' \cite{hennessy2017computer}. Without an intermediary memory subsystem, the processor remains stalled for dozens or hundreds of clock cycles awaiting instruction and operand deliveries.

To bridge this operational frequency discrepancy, memory hierarchies incorporate small, fast SRAM arrays designated as caches. By caching recently referenced memory addresses and their contiguous neighbors, cache controllers exploit temporal locality (recently accessed items are likely to be accessed again) and spatial locality (items with adjacent addresses are likely to be accessed consecutively) \cite{smith1982cache}.

Among the fundamental cache organizations---fully-associative, set-associative, and direct-mapped---the direct-mapped topology provides the absolute lowest access latency and simplest hardware complexity. In a direct-mapped cache, each main memory address maps to exactly one predetermined line in the cache array, entirely eliminating the complex multiplexing logic and replacement policy hardware (e.g., Least Recently Used [LRU]) inherent to associative designs \cite{hill1989evaluating}.

This paper presents the Phase~I architectural formulation, modular Verilog RTL implementation, and comprehensive simulation verification of a 16-bit Direct-Mapped Cache Controller designed for embedded computing environments. The design integrates a synchronous 4-line, 4-byte/line SRAM storage core, a behavioral byte-addressable main memory model with configurable access latency, and an explicit six-state FSM controller enforcing Write-Through and Write-Allocate policies.

\section{Literature Review}
The foundations of cache organization and evaluation metrics were established in seminal work by Alan Jay Smith \cite{smith1982cache}, who comprehensively analyzed miss ratios, line sizes, write policies, and fetch strategies. Smith demonstrated that even diminutive cache capacities drastically curb processor stall cycles by capturing the majority of memory access distributions.

Hill and Smith \cite{hill1989evaluating} evaluated the trade-offs of associativity in CPU caches, proving that while set-associative caches exhibit lower miss rates for heavily conflicting workloads, direct-mapped caches achieve superior cycle times. Because direct-mapped designs perform tag comparison and data selection in parallel without waiting for way-multiplexers, their critical path latency is markedly shorter, often yielding higher overall effective system performance in high-frequency pipelines.

Przybylski \cite{przybylski1990cache} formulated performance-directed approaches for cache and memory hierarchy design, defining analytical models for latency, bandwidth, and pin limitations. Jouppi \cite{jouppi1990improving} focused on mitigating conflict misses in direct-mapped architectures, proposing small auxiliary buffers (victim caches) to store lines evicted by direct-mapped collisions. Kroft \cite{kroft1981lockup} introduced lockup-free cache concepts, demonstrating how non-blocking architectures pipeline concurrent misses.

In educational and synthesizable digital design, the IEEE 1364-2001 standard \cite{ieee2001verilog} serves as the preeminent hardware description medium. The present work builds directly upon these established architectural principles, delivering a clean, synthesizable, cycle-accurate implementation suitable for hardware realization on FPGA platforms.

\section{Background and Working Principles}

\subsection{Cache Organization Parameters}
The specifications for the proposed direct-mapped cache controller are summarized in Table~\ref{tab:specs}.

\begin{table}[htbp]
\caption{Hardware Architectural Specifications}
\label{tab:specs}
\centering
\begin{tabularx}{\columnwidth}{lX}
\toprule
\textbf{Parameter} & \textbf{Hardware Specification} \\
\midrule
Cache Organization & Direct-Mapped (1-Way Associative) \\
Address Bus Width & 16 bits ($A[15:0]$) \\
Data Bus Width (CPU) & 8 bits (Byte-addressable) \\
Number of Cache Lines & 4 lines ($S = 4$) \\
Block / Line Size & 4 bytes ($B = 4$, 32 bits total) \\
Total Cache Data Capacity & 16 bytes ($4 \times 4$ bytes) \\
Tag Field Width & 12 bits ($Tag[11:0] = A[15:4]$) \\
Index Field Width & 2 bits ($Index[1:0] = A[3:2]$) \\
Offset Field Width & 2 bits ($Offset[1:0] = A[1:0]$) \\
Status Bits per Line & 1 Valid bit, 0 Dirty bits \\
Write Policy & Write-Through (Immediate sync) \\
Write Allocation Policy & Write-Allocate (Fetch on write miss) \\
Target HDL Standard & Verilog-2001 (IEEE 1364-2001) \\
\bottomrule
\end{tabularx}
\end{table}

\subsection{Address Decomposition}
For an incoming 16-bit physical memory address $A[15:0]$, the address space is mapped directly into three constituent bit-fields:
\begin{equation}
A[15:0] = \underbrace{A[15:4]}_{\text{Tag (12 bits)}} \;+\; \underbrace{A[3:2]}_{\text{Index (2 bits)}} \;+\; \underbrace{A[1:0]}_{\text{Offset (2 bits)}}
\end{equation}

\begin{enumerate}
    \item \textbf{Byte Offset ($A[1:0]$, 2 bits):} Addresses one of the four individual data bytes ($B_0, B_1, B_2, B_3$) residing inside a 4-byte line block ($\log_2(4) = 2$).
    \item \textbf{Line Index ($A[3:2]$, 2 bits):} Directly decodes which of the four physical cache lines is addressed ($\log_2(4) = 2$).
    \item \textbf{Tag ($A[15:4]$, 12 bits):} Represents the memory block's high-order address page ($16 - 2 - 2 = 12$ bits). The tag is stored in the cache tag memory and compared against incoming addresses to verify hit conditions.
\end{enumerate}

\subsection{Write-Through vs. Write-Back Policy}
Under a \textbf{Write-Through} policy, whenever the processor initiates a memory write, the updated data byte is simultaneously written to both the cache memory and the backing main memory. 
\begin{itemize}
    \item \textit{Advantage:} Main memory always maintains a 100\% coherent, up-to-date representation of all memory locations. Consequently, evicted cache blocks can be overwritten instantaneously without requiring eviction write-backs.
    \item \textit{Hardware Simplification:} The design completely avoids the silicon overhead and verification complexity of tracking ``dirty bits'' per line.
\end{itemize}

\subsection{Write-Allocate Miss Policy}
When a write miss occurs, the controller employs a \textbf{Write-Allocate} (fetch-on-write) strategy. The controller fetches the entire 4-byte block enclosing the referenced address from main memory into the cache line first. Once allocated, the processor's new write byte overwrites the targeted offset in the cache, and the write is concurrently committed to main memory. This ensures subsequent reads and writes to adjacent bytes in the same block immediately achieve cache hits.

\section{Proposed Solution and System Architecture}
The system architecture is structured into a modular hierarchy comprising three distinct hardware modules interconnected via standard request-acknowledge handshake buses:

\begin{enumerate}
    \item \textbf{\texttt{cache\_memory.v}:} A synchronous SRAM module containing four lines of 4-byte data arrays ($4 \times 4 \times 8 = 128$ data flip-flops), four 12-bit tag registers ($4 \times 12 = 48$ bits), and a 4-bit valid vector. It supports single-cycle asynchronous line lookups, full 4-byte block synchronous writes, and single-byte selective updates.
    \item \textbf{\texttt{cache\_controller.v}:} The master digital control engine implementing a six-state Moore/Mealy hybrid FSM. It latches CPU transaction parameters, drives address decoders, performs hit/miss comparisons, commands main memory burst block transfers, and generates single-cycle completion strobes.
    \item \textbf{\texttt{direct\_mapped\_cache.v}:} The structural top-level module encapsulating the cache memory and controller, presenting clean CPU-side and memory-side bus boundaries.
    \item \textbf{\texttt{main\_memory.v}:} A behavioral secondary storage model simulating physical memory access latency (2 clock cycles), providing 4-byte burst read transfers and single-byte write capabilities.
\end{enumerate}

\section{Hardware Design and Verilog Implementation}

\subsection{Finite State Machine (FSM) Engine}
The cache controller transitions across six deterministic states:

\begin{enumerate}
    \item \textbf{\texttt{S\_IDLE} (\texttt{3'b000}):} The controller remains dormant awaiting processor requests (\texttt{cpu\_req == 1'b1}). Upon request arrival, it registers \texttt{cpu\_addr}, \texttt{cpu\_rw}, and \texttt{cpu\_wdata}, advancing to \texttt{S\_LOOKUP}.
    \item \textbf{\texttt{S\_LOOKUP} (\texttt{3'b001}):} The indexed tag and valid bit are evaluated:
    \begin{equation}
        \text{hit} = \text{valid} \land (\text{stored\_tag} == \text{req\_tag})
    \end{equation}
    If $\text{is\_hit}$ is true:
    \begin{itemize}
        \item For Read (\texttt{cpu\_rw == 0}): The addressed byte is retrieved; advances to \texttt{S\_COMPLETE}.
        \item For Write (\texttt{cpu\_rw == 1}): The cache line byte is updated immediately via \texttt{cache\_we\_byte}; advances to \texttt{S\_WRITE\_THROUGH}.
    \end{itemize}
    If $\text{is\_hit}$ is false (Miss): Transitions to \texttt{S\_READ\_MISS}.
    \item \textbf{\texttt{S\_READ\_MISS} (\texttt{3'b010}):} The controller asserts \texttt{mem\_req = 1'b1} and \texttt{mem\_rw = 1'b0}, driving the 4-byte block aligned address $\{Tag, Index, 2'b00\}$. It stalls until main memory signals \texttt{mem\_ready == 1'b1}, capturing the incoming 32-bit burst data.
    \item \textbf{\texttt{S\_ALLOCATE} (\texttt{3'b011}):} The controller writes the 4-byte block, updates the tag register to \texttt{req\_tag}, and asserts the valid bit. If the operation is a write miss, it also modifies the target byte with \texttt{cpu\_wdata} and transitions to \texttt{S\_WRITE\_THROUGH}; if a read miss, it transitions to \texttt{S\_COMPLETE}.
    \item \textbf{\texttt{S\_WRITE\_THROUGH} (\texttt{3'b100}):} Drives single-byte write signals to main memory (\texttt{mem\_req = 1'b1}, \texttt{mem\_rw = 1'b1}, \texttt{mem\_addr = req\_addr}, \texttt{mem\_wdata = req\_wdata}). Upon \texttt{mem\_ready == 1'b1}, memory write is confirmed, advancing to \texttt{S\_COMPLETE}.
    \item \textbf{\texttt{S\_COMPLETE} (\texttt{3'b101}):} Asserts \texttt{cpu\_ready = 1'b1} for one clock cycle, presents \texttt{cpu\_rdata} (for reads), outputs \texttt{cpu\_hit}, and returns to \texttt{S\_IDLE}.
\end{enumerate}

\section{Individual Contributions}
Group Alpha executed this project with responsibilities delineated as follows:

\begin{itemize}
    \item \textbf{Member 1: Palaparthi Harsha Karthikeya (Roll No: 2401010313) --- Lead RTL \& Architecture:}
    Formulated the architectural parameters (16-bit address slicing, 4 lines $\times$ 4 bytes). Designed and implemented synthesizable Verilog modules: \texttt{cache\_memory.v}, \texttt{cache\_controller.v}, and top-level \texttt{direct\_mapped\_cache.v}. Safeguarded legacy 32-bit architectural exploratory assets for Phase~II continuation.
    \item \textbf{Member 2: Dakarapu Hemamdhar nath (Roll No: 2401010137) --- Hardware Verification \& Simulation:}
    Engineered the self-checking automated verification testbench (\texttt{testbench/cache\_tb.v}) covering all twelve functional scenarios. Implemented the cycle-accurate verification engine and VCD waveform generator (\texttt{scripts/simulate\_cache.py}), achieving 100\% assertion coverage.
    \item \textbf{Member 3: Akshay Y (Roll No: 2401020008) --- Memory Subsystem, System Integration \& Documentation:}
    Designed, modeled, and verified the 64~KB secondary memory subsystem (\texttt{main\_memory.v}) with 4-byte burst read packaging, boundary alignment masking, single-byte write-through synchronization, and parameterized 2-cycle latency. Engineered the top-level structural integration (\texttt{direct\_mapped\_cache.v}), unified reset trees, and bus request-acknowledge handshaking. Authored the comprehensive IEEE Phase~I technical report, architectural timing diagrams (\texttt{memory\_interface\_timing.svg}), standalone memory verification suite (\texttt{verify\_memory\_subsystem.py}), 5-minute video demonstration script, and system documentation.
\end{itemize}

\section{Preliminary Results and Discussion}

\subsection{Functional Verification Matrix}
The testbench exercised twenty individual assertions across twelve benchmark scenarios. Table~\ref{tab:results} summarizes the exact simulation results.

\begin{table}[htbp]
\caption{Simulation Verification Results Across 12 Scenarios}
\label{tab:results}
\centering
\begin{tabularx}{\columnwidth}{c X c c c c}
\toprule
\textbf{\#} & \textbf{Scenario Description} & \textbf{Addr} & \textbf{Op} & \textbf{Hit} & \textbf{Status} \\
\midrule
01 & Initial Cold Miss & \texttt{0x0034} & Read  & 0 & \textbf{PASS} \\
02 & Spatial Locality Hit & \texttt{0x0035} & Read  & 1 & \textbf{PASS} \\
03 & Temporal Locality Hit & \texttt{0x0034} & Read  & 1 & \textbf{PASS} \\
04 & Write Hit (Write-Through) & \texttt{0x0035} & Write & 1 & \textbf{PASS} \\
05 & Read Back After Write Hit & \texttt{0x0035} & Read  & 1 & \textbf{PASS} \\
06 & Memory Coherence Assertion & \texttt{0x0035} & Check & 1 & \textbf{PASS} \\
07 & Write Miss (Write-Allocate) & \texttt{0x0028} & Write & 0 & \textbf{PASS} \\
08 & Read Back After Write Miss & \texttt{0x0028} & Read  & 1 & \textbf{PASS} \\
09 & Adjacent Byte Allocation Hit & \texttt{0x0029} & Read  & 1 & \textbf{PASS} \\
10 & Conflict Miss (Tag Mismatch) & \texttt{0x0044} & Read  & 0 & \textbf{PASS} \\
11 & Line Replacement Re-access & \texttt{0x0034} & Read  & 0 & \textbf{PASS} \\
12 & Byte Selection (Offset 1) & \texttt{0x0035} & Read  & 1 & \textbf{PASS} \\
13 & Byte Selection (Offset 2) & \texttt{0x0036} & Read  & 1 & \textbf{PASS} \\
14 & Byte Selection (Offset 3) & \texttt{0x0037} & Read  & 1 & \textbf{PASS} \\
15 & Multi-Line Loading (Line 0) & \texttt{0x0010} & Read  & 0 & \textbf{PASS} \\
16 & Multi-Line Loading (Line 3) & \texttt{0x003C} & Read  & 0 & \textbf{PASS} \\
17 & Consecutive Hit: Line 0 & \texttt{0x0010} & Read  & 1 & \textbf{PASS} \\
18 & Consecutive Hit: Line 1 & \texttt{0x0034} & Read  & 1 & \textbf{PASS} \\
19 & Consecutive Hit: Line 2 & \texttt{0x0028} & Read  & 1 & \textbf{PASS} \\
20 & Consecutive Hit: Line 3 & \texttt{0x003C} & Read  & 1 & \textbf{PASS} \\
\bottomrule
\end{tabularx}
\end{table}

\subsection{Conflict Miss \& Line Replacement Analysis}
A pivotal architectural characteristic of direct-mapped caches is vulnerability to conflict misses when two active memory addresses map to the identical index. This was demonstrated using addresses \texttt{0x0034} and \texttt{0x0044}:
\begin{itemize}
    \item \texttt{0x0034}: Binary \texttt{0000 0000 0011 0100} $\rightarrow$ Tag = \texttt{12'h003}, Index = \texttt{2'b01} (Line 1), Offset = \texttt{2'b00}.
    \item \texttt{0x0044}: Binary \texttt{0000 0000 0100 0100} $\rightarrow$ Tag = \texttt{12'h004}, Index = \texttt{2'b01} (Line 1), Offset = \texttt{2'b00}.
\end{itemize}

Because both addresses possess index bits \texttt{2'b01}, they compete for Line~1. In Test~01, block \texttt{0x0034--0x0037} was loaded into Line~1. In Test~10, an access to \texttt{0x0044} caused an immediate tag mismatch against stored tag \texttt{12'h003}. Because the line was valid but the tag mismatched, a conflict miss was flagged, evicting the existing block and allocating \texttt{0x0044--0x0047}. Subsequently, re-accessing \texttt{0x0034} in Test~11 produced another miss, confirming the eviction and replacement sequence.

\subsection{Performance and Latency Analysis}
\begin{itemize}
    \item \textbf{Read Hit Latency:} 2 clock cycles (\texttt{S\_IDLE} $\rightarrow$ \texttt{S\_LOOKUP} $\rightarrow$ \texttt{S\_COMPLETE}).
    \item \textbf{Read Miss Latency:} 5 clock cycles (\texttt{S\_IDLE} $\rightarrow$ \texttt{S\_LOOKUP} $\rightarrow$ \texttt{S\_READ\_MISS} [2 memory cycles] $\rightarrow$ \texttt{S\_ALLOCATE} $\rightarrow$ \texttt{S\_COMPLETE}).
    \item \textbf{Write Hit Latency:} 4 clock cycles (\texttt{S\_IDLE} $\rightarrow$ \texttt{S\_LOOKUP} $\rightarrow$ \texttt{S\_WRITE\_THROUGH} [2 memory cycles] $\rightarrow$ \texttt{S\_COMPLETE}).
    \item \textbf{Write Miss Latency:} 7 clock cycles (\texttt{S\_IDLE} $\rightarrow$ \texttt{S\_LOOKUP} $\rightarrow$ \texttt{S\_READ\_MISS} [2 memory cycles] $\rightarrow$ \texttt{S\_ALLOCATE} $\rightarrow$ \texttt{S\_WRITE\_THROUGH} [2 memory cycles] $\rightarrow$ \texttt{S\_COMPLETE}).
\end{itemize}

\section{Limitations and Phase II Continuation}
While the Phase~I direct-mapped architecture provides minimal access latency and zero dirty-bit overhead, several inherent architectural limitations exist:
\begin{enumerate}
    \item \textbf{Conflict Miss Thrashing:} Alternating accesses between addresses sharing the same index (e.g., \texttt{0x0034} and \texttt{0x0044}) causes recurring conflict misses.
    \item \textbf{Write-Through Bus Traffic:} Every write operation traverses the memory bus, creating bus contention in multi-core or bandwidth-constrained systems.
\end{enumerate}

\textbf{Phase II Continuation Roadmap:}
In Phase~II, the architecture will be expanded to a 2-way set-associative cache controller featuring a Write-Back policy with dirty bits and Least Recently Used (LRU) replacement, comparing hit rates and bus transactions against this Phase~I baseline.

\section{Conclusion}
This paper presented the design, Verilog-2001 RTL implementation, and cycle-accurate verification of a 16-bit direct-mapped cache controller. The modular design incorporates a synchronous 4-line SRAM array, an explicit six-state FSM, and a latency-modeled secondary memory model. Functional verification proved 100\% assertion success across twenty verification points, validating cold misses, spatial locality hits, write-through memory synchronization, and direct-mapped conflict replacements. The synthesized RTL provides a verified baseline ready for Phase~II associative expansion.

\bibliographystyle{IEEEtran}
\bibliography{references}

\end{document}
